% SPDX-License-Identifier: CC-BY-4.0
% Original manuscript by Chatgpt. License: https://creativecommons.org/licenses/by/4.0/
% Project: https://github.com/rongyulin/physics-open-problems-AI/tree/main/problems/local-sic-dimension-12
\documentclass[11pt]{article}
\usepackage[T1]{fontenc}
\usepackage{lmodern,microtype,amsmath,amssymb,amsthm,mathtools,booktabs}
\usepackage[a4paper,top=24mm,bottom=25mm,left=27mm,right=27mm]{geometry}
\usepackage{xcolor}
\usepackage[colorlinks=true,linkcolor=blue!45!black,citecolor=blue!45!black,urlcolor=blue!45!black]{hyperref}
\hypersetup{pdftitle={Abelian covariance of symmetric informationally complete measurements in dimension twelve},pdfauthor={Chatgpt}}
\newtheorem{theorem}{Theorem}
\newtheorem{corollary}[theorem]{Corollary}
\newtheorem{proposition}[theorem]{Proposition}
\newcommand{\CC}{\mathbb C}\newcommand{\ZZ}{\mathbb Z}\newcommand{\RR}{\mathbb R}
\newcommand{\tr}{\operatorname{tr}}\newcommand{\rank}{\operatorname{rank}}
\newcommand{\cof}{\operatorname{cof}}\newcommand{\spec}{\operatorname{spec}}
\newcommand{\cP}{\mathcal P}\newcommand{\one}{\mathbf 1}
\title{\bfseries Abelian covariance of symmetric informationally complete measurements in dimension twelve}
\author{Chatgpt}
\date{27 November 2026}

% Place the license below the first-page folio without moving the article.
\AddToHookNext{shipout/foreground}{%
  \put(0,\dimexpr 19pt-\paperheight\relax){%
    \makebox[\paperwidth][c]{\normalfont\footnotesize
      Paper and \LaTeX{} source licensed under \href{https://creativecommons.org/licenses/by/4.0/}{CC BY 4.0}%
      \enspace\textperiodcentered\enspace
      \href{https://github.com/rongyulin/physics-open-problems-AI/tree/main/problems/local-sic-dimension-12}{rongyulin/physics-open-problems-AI}}}}
\begin{document}
\maketitle
\begin{abstract}
Can two qubits and a qutrit generate a symmetric informationally complete measurement by local Weyl--Heisenberg displacements? We prove that they cannot. Such a fiducial would have a two-qubit marginal of rank at most three with all fifteen nontrivial Pauli expectations of modulus $1/\sqrt{13}$. A zero eigenvalue would force an eigenvalue $1/2$ by algebraic conjugation, incompatible with the marginal purity $7/13$. This excludes one of the two possible abelian transitive projective symmetry groups in dimension twelve. The other, $\mathbb Z_{12}^2$, occurs in known exact constructions. We also determine the exact least squared deviation of the fifteen local overlap moduli from the SIC value, derive a dimension-dependent obstruction, and give an explicit gap below the local stabilizer R\'enyi entropy bound of $\log_2(13/2)$. Lean~4 verifies the nonexistence theorem, the attained global minimum, and the entropy bound directly from normalized state amplitudes.
\end{abstract}

\section{Two symmetries on the same Hilbert space}

A symmetric informationally complete measurement (SIC) in dimension $d$ consists of $d^2$ rank-one effects $\Pi_j/d$, where $\Pi_j=|\psi_j\rangle\langle\psi_j|$ and the unit vectors have pairwise squared overlaps $1/(d+1)$~\cite{renes}. These projectors form a basis for the operator space: their Hilbert--Schmidt Gram matrix has diagonal $1$ and off-diagonal $1/(d+1)$, and is invertible. Also $\sum_j\Pi_j=dI$, as can be seen by pairing both sides with every $\Pi_i$. Thus the $d^2$ probabilities determine a quantum state. The existence problem in arbitrary dimension is prominent in quantum information~\cite{open}.

In dimension twelve, exact SICs generated by a cyclic shift and clock on twelve levels are known~\cite{appleby}. A different experiment is available on $\CC^2\otimes\CC^2\otimes\CC^3$. One can act on each of the three factors with its own shift and clock. The two constructions have the same number of displacements, but their projective index groups differ:
\[
     \ZZ_{12}^{\,2}
     \qquad\text{and}\qquad
     \ZZ_2^{\,4}\times\ZZ_3^{\,2}\simeq(\ZZ_2\times\ZZ_6)^2.
\]
For example, the largest order of an element in the second group is six. Banacki et al.~\cite[Sec.~5.1]{banacki} explicitly asked whether the latter action admits a SIC when studying the spin state in Higgs decay. An earlier variational search by Chien and Waldron~\cite[Sec.~5, Table~2]{chien} includes the local group as \texttt{SmallGroup(144,197)} and gives no numerical SIC for it. Table~2 does record one for the usual displacement group \texttt{SmallGroup(144,101)}. Khodaiemehr subsequently reported another unsuccessful numerical search for the same group~\cite[Sec.~III]{ghost}. The distinction between ordinary and product displacements also enters the study of extremal non-stabilizerness by Damski et al.~\cite{damski}. Our argument settles the local existence question for arbitrary pure states of the three factors.

\section{The obstruction in fifteen expectations}

Write $\cP_2=\{I,X,Y,Z\}^{\otimes 2}$ for the sixteen Hermitian two-qubit Pauli operators, with global phases removed. On the qutrit let $X_3|j\rangle=|j+1\bmod3\rangle$ and $Z_3|j\rangle=e^{2\pi i j/3}|j\rangle$. A local orbit consists of the 144 rays
\[
       (P\otimes X_3^a Z_3^b)\psi,\qquad P\in\cP_2,\quad a,b\in\ZZ_3.
\]
The operator order and scalar phase conventions for the qutrit displacements do not affect the rays.

\begin{theorem}\label{th:no}
There is no pure state on $\CC^4\otimes\CC^3$ satisfying
\begin{equation}
  |\langle\psi|(P\otimes I_3)|\psi\rangle|^2=1/13
       \quad\text{for all } P\in\cP_2\setminus\{I\}.              \label{eq:flat}
\end{equation}
In particular the orbit of the local Weyl--Heisenberg group on two qubits and a qutrit cannot be a SIC.
\end{theorem}
\begin{proof}
Let $\rho$ be the two-qubit reduced density matrix. If $F$ is the $4$-by-$3$ matrix of amplitudes of $\psi$, then $\rho=FF^\dagger$, so it is positive semidefinite, has trace one and has rank at most three. Orthogonality of the Pauli basis, $\tr(PQ)=4\delta_{PQ}$, and the reality of $\tr(P\rho)$ give
\begin{equation}
 \rho=\frac14\left(I+\frac{S}{\sqrt{13}}\right),
 \qquad S=\sum_{P\ne I}s_P P,\qquad s_P\in\{-1,1\}.
 \label{eq:rho}
\end{equation}
Hence $\tr S=0$, $\tr S^2=60$ and
\begin{equation}
                    \tr\rho^2=\frac{1+15/13}{4}=\frac7{13}.
                    \label{eq:purity}
\end{equation}

Singularity of $\rho$ makes $-\sqrt{13}$ an eigenvalue of $S$. All entries of $S$ are Gaussian integers. Its characteristic polynomial has coefficients in $\ZZ[i]$, and those coefficients are real by Hermiticity. They belong to $\ZZ$. Thus $x^2-13$ divides that polynomial over $\mathbb Q$, giving the conjugate eigenvalue $+\sqrt{13}$. By Eq.~\eqref{eq:rho}, $\rho$ has an eigenvalue $1/2$. The remaining nonnegative eigenvalues sum to $1/2$, so their squares sum to at most $1/4$. This gives $\tr\rho^2\leq1/2$, contradicting Eq.~\eqref{eq:purity}. The local SIC condition implies Eq.~\eqref{eq:flat} by comparing its fiducial with the fifteen displacements acting only on the qubits.
\end{proof}

For pure states on the qubits alone, Godsil and Roy obtained an earlier Pauli-orbit obstruction~\cite[Lemma~3.1]{godsil}. Here tracing out the qutrit permits a mixed marginal of rank up to three. The rank obstruction applies to the whole twelve-dimensional Hilbert space. Restrictions from a particular Hamiltonian, decay amplitude or entanglement class do not enter the proof.

\subsection{A determinant argument and sharpness for the local constraints}

For any choices of the signs in Eq.~\eqref{eq:rho}, $\tr S^3$ is a nonzero odd multiple of 24. To see this, label the nonidentity Pauli operators by the 15 nonzero vectors of $\mathbb F_2^4$, with commutation given by the standard symplectic form. Only triples with labels $u,v,u+v$ can contribute to the trace. For anticommuting pairs their contributions cancel between opposite orderings. For a commuting triple its six orderings contribute the same $\pm4$. Each label commutes with six other nonzero, distinct labels, so the number of unordered commuting triples is $15\cdot6/(3\cdot2)=15$. A sum of fifteen signs is odd.

Consequently a signed matrix $S$ of this form cannot have an eigenvalue $\sqrt{k}$ or $-\sqrt{k}$ for any positive nonsquare integer $k$. Integer coefficients would give both roots. The two remaining eigenvalues sum to zero, which would force $\tr S^3=0$. This proof of Theorem~\ref{th:no} uses only the singularity of the proposed marginal.

\begin{proposition}\label{prop:sharp}
For an integer $q\geq1$, there is a pure state on $\CC^4\otimes\CC^q$ with all fifteen nonidentity two-qubit Pauli expectations of modulus $1/\sqrt{4q+1}$ if and only if $q=2$ or $q\geq4$.
\end{proposition}
\begin{proof}
For $q=1,3$ the reduced state is singular and $4q+1$ is respectively 5 or 13. The determinant obstruction applies. Conversely take all fifteen signs positive. Writing $T=I+X+Y+Z$, one gets $S=T\otimes T-I$. Its eigenvalues are $3+2\sqrt3,\ 3-2\sqrt3,\ -3,\ -3$, since $X+Y+Z$ has eigenvalues $\pm\sqrt3$. The matrix $(I+S/\sqrt{4q+1})/4$ is a trace-one positive matrix of rank two at $q=2$ and of rank four at $q\geq4$. It admits a purification in the claimed space. This realizes precisely the fifteen specified expectations.
\end{proof}

For $q=3$ the fifteen conditions in Eq.~\eqref{eq:flat} are also minimal under deletion. To omit a specified Pauli $P_0$, start at $R_+=(I+S/\sqrt{13})/4$ with all signs positive. It is positive definite by the spectrum above. The matrix $R_++tP_0$ keeps trace one and preserves the other fourteen expectations. It has a negative eigenvalue for sufficiently large positive $t$, since $P_0$ has a negative eigenvalue. Its first point on the boundary of the positive cone has rank at most three and can be purified with a qutrit. Thus any prescribed fourteen conditions are simultaneously attainable.

\section{Abelian covariance in dimension twelve}

\begin{corollary}\label{cor:abel}
Every abelian subgroup of the projective unitary group acting transitively on a twelve-dimensional SIC is isomorphic to $\ZZ_{12}^2$ and, up to a change of basis and relabelling, acts by the usual twelve-level Weyl displacements. Such a group occurs.
\end{corollary}
\begin{proof}
An element fixing each SIC projector fixes a basis of the operator space and hence is projectively the identity. An abelian transitive permutation group acting faithfully is regular: a stabilizer fixes the entire orbit. Its order is therefore 144. For $g\ne e$, conjugation by a unitary representative $U_g$ permutes the projector basis without fixed points. The trace of this conjugation map is $|\tr U_g|^2$, hence $\tr U_g=0$. It follows that the 144 representatives form an orthogonal basis of the matrix algebra.

For clarity we recall the abelian error-basis classification~\cite{kr}. Projective commutators give an alternating bicharacter $\beta(g,h)$ with values on the unit circle. It is nondegenerate: a representative commuting with the entire matrix basis is scalar. Different Sylow subgroups pair trivially. In a finite abelian primary group choose $g$ of largest order $r$. The character $\beta(g,\cdot)$ has order $r$, so some $h$ pairs with $g$ in a primitive $r$th root of unity. Then $\langle g,h\rangle\simeq\ZZ_r^2$ and is nondegenerate. Multiplication by appropriate powers of $g,h$ moves each element into the orthogonal complement of this pair. Iteration splits the group into such pairs.

Here the 3-primary part has order 9 and gives one pair of order three. The 2-primary part has order 16 and gives either one pair of order four or two pairs of order two. Each pair acts as a clock and shift on its own factor. Indeed, rephase its representatives to have the stated orders and choose generators with standard commutator. Take a simultaneous eigenvector for the commuting clocks. Applying every product of the shifts produces $12$ vectors with distinct joint clock eigenvalues and hence an orthogonal basis. In this basis the action has the asserted tensor form. For the pair sizes $4,3$, the Chinese remainder isomorphism identifies the resulting displacement group with the usual $\ZZ_{12}^2$ action. The case $2,2,3$ is excluded by Theorem~\ref{th:no}. The exact twelve-dimensional construction in Ref.~\cite{appleby} realizes the remaining action.
\end{proof}
This classifies transitive \emph{abelian projective unitary} symmetry groups. The full symmetry group of an individual SIC may contain further symmetries, including noncommuting ones.

\section{More qubits and a smaller ancillary system}

Let $m=2^\ell$, $\ell\ge2$, and let a pure state on $\CC^m\otimes\CC^q$ have all $m^2-1$ nonidentity Pauli expectations on the first factor of modulus $1/\sqrt{\Delta}$, where $\Delta=mq+1$. Its reduced state is $\rho=(I+S/\sqrt{\Delta})/m$, with an integer characteristic polynomial for the signed Pauli sum $S$ as above.

\begin{proposition}\label{prop:family}
When $\Delta=mq+1$ is nonsquare, flat Pauli expectations on the $m$-dimensional factor require
\[
                         q>\frac{m+\sqrt m}{2}.
\]
\end{proposition}
\begin{proof}
For $q\ge m$ this is immediate. Otherwise the eigenvalue $-\sqrt{\Delta}$ of $S$ has multiplicity at least $m-q$. Its conjugate $+\sqrt{\Delta}$ has the same multiplicity, since the characteristic polynomial has rational coefficients. Remove $m-q$ copies of each. There are $r=2q-m$ remaining eigenvalues, so $0\le r<m$.

In $\rho$ the removed eigenvalues correspond to $m-q$ zeros and $m-q$ copies of $2/m$. The remaining eigenvalues are nonnegative and sum to $r/m$. Their squared sum is at most $(r/m)^2$. On the other hand, Pauli orthogonality gives the exact purity $\tr\rho^2=(1+(m^2-1)/\Delta)/m=(m+q)/(mq+1)$. It follows that
\[
 \frac{m+q}{mq+1}\leq\frac{4(m-q)+r^2}{m^2}.
\]
After clearing denominators and substituting $2q=m+r$, twice the difference of the right and left numerators is
\[
                 (m+r-2)\bigl(mr^2+2r-m^2-2m\bigr)\geq0 .
\]
The first factor is positive. Hence $r^2\ge m+2-2r/m>m$, proving the result.
\end{proof}
An infinite family occurs at odd $\ell\ge3$ and $q=m/2+1$, where $\Delta\equiv2\bmod3$ is nonsquare. The first ambient dimensions are $40$, $544$ and $8320$. The proposition excludes local displacement SICs throughout this family. The same bound also excludes, for example, $m=16$, $q=10$, $\Delta=161$. The nonsquare condition is essential: the three-qubit Hoggar SIC occurs at $m=4$, $q=2$ and $\Delta=9$~\cite{review}.

\section{A quantitative gap in stabilizer entropy}
The second-order stabilizer R\'enyi entropy for the specified local displacement basis $\mathcal D$ is~\cite{banacki,cuffaro}
\[
 M_2(\psi)=-\log_2\left(\frac1{12}\sum_{D\in\mathcal D} q_D^2\right),
 \qquad q_D=|\langle\psi|D|\psi\rangle|^2 .
\]
Orthogonality of the displacement basis gives $\sum_Dq_D=12$ for a pure state, with $q_I=1$. Consequently
\begin{equation}
    \sum_D q_D^2=\frac{24}{13}+\sum_{D\ne I}(q_D-1/13)^2 . \label{eq:deficit}
\end{equation}
Compactness and Theorem~\ref{th:no} already give a strict gap below $\log_2(13/2)$. We give an explicit value.

\begin{proposition}\label{prop:gap}
For $r_P=\langle\psi|(P\otimes I_3)|\psi\rangle$, the exact minimum over unit vectors on $\CC^4\otimes\CC^3$ is
\begin{equation}
 \min_{\psi}\sum_{P\ne I}\left(|r_P|-1/\sqrt{13}\right)^2
             =\frac{(1-3/\sqrt{13})^2}{3}. \label{eq:min}
\end{equation}
Put $t=1/\sqrt{13}$ and $h=(1-3t)/\sqrt3$. Every such state satisfies
\[
 M_2(\psi)\leq
 \log_2(13/2)-\log_2\left(1+\frac{13}{24}h^2(2t-h)^2\right)
 \simeq 2.698901070\ {\rm bits}.
\]
\end{proposition}
\begin{proof}
We first show that every signed two-qubit Pauli sum $S$ satisfies
\begin{equation}
    \min_{\lambda\in\spec(S)}|\lambda+\sqrt{13}|\ \geq\ \sqrt{13}-3. \label{eq:gap}
\end{equation}
Here are details of a finite integer certificate. Write $\sigma=(X,Y,Z)$ and
$S=(a\cdot\sigma)\otimes I+I\otimes(b\cdot\sigma)+\sum_{ij}C_{ij}\sigma_i\otimes\sigma_j$
with $a,b\in\{-1,1\}^3$ and $C$ a sign matrix. The characteristic polynomial has the form $x^4-30x^2+B x+D_0$. Set $C'_{ij}=a_iC_{ij}b_j$, $\eta=\prod_i a_i b_i$, and $J=\one\one^{\mathsf T}$ of size three. Its coefficients are
\begin{align}
 B&=-8\left(\sum_{ij}C'_{ij}-\eta\det C'\right), \notag\\
 D_0&=225-4\left(\|C'\one\|^2+\|(C')^{\mathsf T}\one\|^2+
                              \|J-\eta\cof C'\|_F^2\right).              \label{eq:coeff}
\end{align}
For verification, the Pauli expansion of $S^2$ has scalar term $15$, one-body terms $2Cb$, $2C^{\mathsf T}a$, and two-body coefficient matrix $2(ab^{\mathsf T}-\cof C)$. Evaluating $\tr S^3$ and $\tr S^4$ by orthogonality and using Newton's identities gives Eq.~\eqref{eq:coeff}.

As $C'$ and $\eta$ take their $2^{10}$ possible values, the entire list of pairs $(|B|,D_0)$ is
\[
 (72,-27),\quad(56,-27),\quad(40,5),\quad(24,69),\quad(8,37),\quad(8,165).
\]
For each of the twelve signed quartics $f$, the coefficients of
$(4+4u)^4 f((12+17u)/(4+4u))$ have a common weak sign and its leading coefficient is nonzero. Hence $f$ has no zero in $(3,17/4)$. Reversing $B$ gives the same assertion for $-\lambda$. Since $3<\sqrt{13}$ and $17/4-\sqrt{13}>\sqrt{13}-3$, this proves Eq.~\eqref{eq:gap}. The two finite checks are included in the Lean project.

Let $t=1/\sqrt{13}$ and choose $s_P=\operatorname{sign}(r_P)$, taking $+1$ at zero. The trace-zero Hermitian matrix $K=4\sqrt{13}\rho-(\sqrt{13}I+S)$ satisfies $\|K\|_{\mathrm{op}}^2\leq(3/4)\|K\|_F^2$: if $\lambda$ has largest absolute value, the other three eigenvalues sum to $-\lambda$ and their squared sum is at least $\lambda^2/3$. For a unit vector $y\in\ker\rho$, Eq.~\eqref{eq:gap} yields
\[
 (\sqrt{13}-3)^2\leq\|Ky\|^2\leq\tfrac34\|K\|_F^2
       =3\sum_{P\ne I}(\sqrt{13}|r_P|-1)^2.
\]
This proves the lower bound in Eq.~\eqref{eq:min}.

Equality is attained. Use the all-positive sum $S$ from Proposition~\ref{prop:sharp}; the least eigenvalue of $R_+=(I+tS)/4$ is $a=(1-3t)/4$. The Bell singlet $u=(|01\rangle-|10\rangle)/\sqrt2$ lies in this eigenspace because $(T\otimes T)u=(\det T)u=-2u$. Put
\[
       R=R_++(a/3)I-(4a/3)|u\rangle\langle u|.
\]
Its eigenvalue on $u$ is zero. Each of the other three eigenvalues is positive, and $\tr R=1$. It therefore has a qutrit purification. By the singlet projector identity
$4|u\rangle\langle u|=I-X\otimes X-Y\otimes Y-Z\otimes Z$, the three expectations for $X\otimes X,Y\otimes Y,Z\otimes Z$ equal $1/3$, and the other twelve equal $t$. Their squared deviations sum to $3(1/3-t)^2$, as claimed.

Finally let $z_P=|r_P|$, so $q_{P\otimes I}=z_P^2$ and $0<h<t$. If some $|z_P-t|\ge h$, then $|z_P^2-t^2|\ge h(2t-h)$, by considering $z_P\le t-h$ and $z_P\ge t+h$. Otherwise all $z_P+t\ge2t-h$. In that case Eq.~\eqref{eq:min} gives $\sum_{P\ne I}(z_P^2-t^2)^2\ge(2t-h)^2\sum_{P\ne I}(z_P-t)^2\ge(2t-h)^2h^2$. In either case the qubit-only terms already contribute the required deficit to Eq.~\eqref{eq:deficit}.
\end{proof}
This is a sufficient gap. The maximizing states and the exact maximum of $M_2$ for this local displacement group remain to be determined.

\section{Formal verification and reproducibility}
The accompanying project pins Lean 4.28.0 and Mathlib. Its main theorem,
\begin{center}\texttt{SicNoGo.no\_local\_hw\_sic\_dim12},\end{center}
states that no $4$-by-$3$ matrix of complex amplitudes generates a local SIC. The displacement definitions use the three Pauli matrices, the cyclic qutrit shift and the primitive phase $e^{2\pi i/3}$. The predicate imposes normalization and equal squared overlap on every pair of distinct labels of the full 144-vector orbit. The proof derives Eq.~\eqref{eq:flat}, constructs the marginal $FF^\dagger$, proves that its determinant vanishes, and reconstructs the signed Pauli expansion from the measured expectations. It proves the contradictory nonzero determinant by an exact polynomial identity and an odd linear coefficient. No numerical spectrum, optimizer or rank assumption is imported. The project also proves invariance under arbitrary choices of scalar unit phases for the representatives.

The theorem \texttt{SicStability.sharp\_minimum} verifies Eq.~\eqref{eq:min} as an attained global minimum, with an explicit normalized state attaining it. Computing all 144 overlaps gives the entropy bound in \texttt{SicEntropy.}\allowbreak\texttt{local\_entropy\_upper}. Both quantify over arbitrary normalized complex amplitude matrices. Lean also verifies the arithmetic implication from the paired marginal spectrum in Proposition~\ref{prop:family}. The group classification, the general conjugate-eigenvalue argument and the results about deleting a constraint and varying the ancilla dimension are given as mathematical proofs in the text. Instructions for compiling from source and replaying the declarations with \texttt{leanchecker} accompany the project.

\begingroup\small
\begin{thebibliography}{99}
\setlength{\itemsep}{3pt}
\bibitem{renes} J. M. Renes, R. Blume-Kohout, A. J. Scott and C. M. Caves, Symmetric informationally complete quantum measurements, \emph{J. Math. Phys.} \textbf{45}, 2171--2180 (2004). \href{https://arxiv.org/abs/quant-ph/0310075}{arXiv:quant-ph/0310075}.
\bibitem{open} P. Horodecki, \L. Rudnicki and K. \.Zyczkowski, Five open problems in quantum information theory, \emph{PRX Quantum} \textbf{3}, 010101 (2022). \href{https://doi.org/10.1103/PRXQuantum.3.010101}{doi:10.1103/PRXQuantum.3.010101}.
\bibitem{appleby} D. M. Appleby et al., Systems of imprimitivity for the Clifford group, \href{https://arxiv.org/abs/1210.1055}{arXiv:1210.1055} (2012).
\bibitem{banacki} M. Banacki et al., Qubit-qubit-qutrit quantum correlations in $H\to f\bar f V$, \emph{Phys. Rev. D} \textbf{114}, 093007 (2026). \href{https://doi.org/10.1103/hv6p-c34d}{doi:10.1103/hv6p-c34d}.
\bibitem{chien} T.-Y. Chien and S. Waldron, Nice error frames, canonical abstract error groups and the construction of SICs, \emph{Linear Algebra Appl.} \textbf{516}, 93--117 (2017). \href{https://doi.org/10.1016/j.laa.2016.11.026}{doi:10.1016/j.laa.2016.11.026}.
\bibitem{ghost} H. Khodaiemehr, From ghost SICs to SIC-POVMs: A Galois reduction of Zauner's conjecture, \href{https://arxiv.org/abs/2611.01524}{arXiv:2611.01524} (2026).
\bibitem{kr} A. Klappenecker and M. R\"otteler, Beyond stabilizer codes I: Nice error bases, \emph{IEEE Trans. Inform. Theory} \textbf{48}, 2392--2395 (2002). \href{https://arxiv.org/abs/quant-ph/0010082}{arXiv:quant-ph/0010082}.
\bibitem{godsil} C. Godsil and A. Roy, Equiangular lines, mutually unbiased bases, and spin models, \href{https://arxiv.org/abs/quant-ph/0511004}{arXiv:quant-ph/0511004} (2005).
\bibitem{review} C. A. Fuchs, M. C. Hoang and B. C. Stacey, The SIC question: History and state of play, \emph{Axioms} \textbf{6}, 21 (2017). \href{https://doi.org/10.3390/axioms6030021}{doi:10.3390/axioms6030021}.
\bibitem{cuffaro} G. Cuffaro and C. A. Fuchs, Quantum states with maximal magic, \href{https://arxiv.org/abs/2412.21083}{arXiv:2412.21083} (2024).
\bibitem{damski} B. S. Damski, R. Bistro\'n, D. Ponterio, J. Czartowski and K. \.Zyczkowski, Product Weyl--Heisenberg covariant mutually unbiased bases and extremal non-stabilizerness, \href{https://arxiv.org/abs/2603.15550}{arXiv:2603.15550} (2026).
\end{thebibliography}
\endgroup
\end{document}
